% ====================================================================
% FILE: content/09_thresholds_ru.tex
% Расширенный ландшафт порогов — Core 1.3
% ====================================================================

\section{Расширенный ландшафт порогов}
\label{sec:thresholds-extended}

Эта глава восстанавливает полную теорию порогов в Онтологии континуумов (OC),
консолидируя структурные элементы из направлений Физики, химии,
Биологии и когнитивных направлений. Пороги — это центральный механизм, через который
порождаются качественные изменения, появление новых размерностей, коллапс и
смерть континуумов.

Пороги универсальны: они применяются к физическим, химическим, биологическим,
когнитивным, социальным и теоретическим континуумам на всей иерархии
\(K_0\)–\(K_{10}\).

\subsection{Определение порога}

\emph{Порог} представлен функцией ограничения
\[
   f_k : \Omega(K) \rightarrow \mathbb{R}
\]
с условием
\[
   f_k(s) \le 0
   \quad\text{для всех допустимых состояний } s\in\Omega(K).
\]

Пороги разделяют качественно разные динамические режимы:
\begin{itemize}
    \item \( f_k(s) < 0 \): безопасный режим (внутренняя область),
    \item \( f_k(s) = 0 \): пограничный режим,
    \item \( f_k(s) > 0 \): запретный режим (вне \(\Omega(K)\)).
\end{itemize}

Следовательно, пороги совместно определяют:
\begin{itemize}
    \item форму допустимого множества \(\Omega(K)\),
    \item границу \(\partial\Omega(K)\),
    \item условия, при которых континуум существует или коллапсирует.
\end{itemize}

Пороги не являются статическими константами.
Они зависят от потенциалов \(P(t)\), потоков \(J(t)\), осей \(A(K)\),
экологических ограничений и структуры пространства вложения \(M\).
На структурном уровне это кодируется оператором эволюции порога
\(H\) из общего аппарата эволюции:
\[
  \frac{d\Theta_k}{dt}
   = H_k\big(\Theta_k, P, J, A, M\big),
\]
где \(\Theta_k\) обозначает набор параметров, определяющий ограничение \(f_k\).

\subsection{Таксономия порогов}

OC выделяет семь универсальных классов порогов. Эта таксономия была
сформирована в Core~2.x и уточнена по доменным запускам.

\paragraph{1. Пороги существования \(\Theta_{\rm exist}\).}
Условия, необходимые для того, чтобы \(\Omega(K)\neq\emptyset\).
Примеры:
\begin{itemize}
    \item минимальные энергетические условия для связных физических состояний (\(K_2\)),
    \item минимальное каталитическое замыкание (существование RAF) в \(K_3\),
    \item минимальная непрерывность и целостность мембраны в \(K_4\),
    \item минимальная когерентность внутренних моделей в \(K_6\),
    \item минимальная легитимность и доверие для институтов в \(K_7\).
\end{itemize}

\paragraph{2. Пороги устойчивости \(\Theta_{\rm stab}\).}
Условия, при которых потоки остаются ограниченными и циклы устойчивы.
Примеры:
\begin{itemize}
    \item условия конфайнмента и ограниченности в физике (\(K_2\)),
    \item условия осмотического гомеостаза в протоклетках (\(K_4\)),
    \item устойчивость мембранного потенциала в нейронах (\(K_5\)),
    \item нормативная устойчивость в социальных континуумах (\(K_7\)).
\end{itemize}

\paragraph{3. Критические пороги \(\Theta_{\rm crit}\).}
Поверхности качественного изменения при фиксированной размерности.
Примеры:
\begin{itemize}
    \item фазовые переходы и пороги типа BKT (\(K_2\)),
    \item запуск каталитических сетей (\(K_3\)),
    \item метаболическое переключение в \(K_4\),
    \item порог инициации спайка в \(K_5\),
    \item бифуркации ошибки прогноза в \(K_6\),
    \item институциональные бифуркации в \(K_7\).
\end{itemize}

\paragraph{4. Размерностные пороги \(\Theta_{\rm dim}\).}
Ограничения, которые запускают создание новых осей.
Размерностные пороги центральны в OC:
\[
   T(K,t) > \Theta_{\rm dim}(K)
   \quad\Rightarrow\quad
   \text{возникновение новой оси } A_{\rm new}.
\]

Они определяют переходы \(K_x \to K_{x+1}\) на всей иерархии.

\paragraph{5. Пороги смерти \(\Theta_{\rm death}\).}
Значения, за которыми не остается допустимых состояний. Эквивалентно,
\[
   \forall s\in\overline{\Omega(K)}:\ \exists k\ \text{с таким, что } f_k(s) > 0
   \quad\Longleftrightarrow\quad
   \Omega(K)=\emptyset.
\]

\paragraph{6. Пороги экспрессивности \(\Theta_{\rm expr}\).}
Пороги на репрезентативную емкость. Континуум коллапсирует, когда
\[
   \dim\big(A(K)\big)
   <
   \dim\big({\rm Differences}(K)\big),
\]
где \(\mathrm{Differences}(K)\) обозначает эффективную размерность
структурно релевантных различий, накладываемых пространством вложения и
ограничениями. Этот класс критичен для когнитивного, социального и
мета-теоретического коллапса.

\paragraph{7. Пороги вложения \(\Theta_{\rm embed}\).}
Ограничения, накладываемые пространством вложения \(M_x\).
Континуум не может существовать как живой континуум, если пространство вложения
не поддерживает его оси, пороги или потоки; формально это отражено в условиях
согласованности в Теоремах~4 и~5 (раздел Результаты).

\subsection{Геометрия порогов}

Пороги индуцируют стратифицированную геометрию на \(\Omega(K)\).
Для каждой пороговой функции \(f_k\) мы рассматриваем множества уровней
\[
  \Sigma_k(c) = \{\, s \in \Omega(K) \mid f_k(s) = c \,\}.
\]

В частности:
\begin{itemize}
    \item \(\Sigma_k(0)\) задает компоненты границы \(\partial\Omega(K)\),
    \item \(\Sigma_k(c<0)\) задает безопасные области,
    \item \(\Sigma_k(c>0)\) задает запретные области.
\end{itemize}

Взаимодействие порогов порождает:
\begin{itemize}
    \item много-пороговые пересечения,
    \item кауспидоподобные геометрии (катастрофоподобные формы),
    \item области сконцентрированной кривизны,
    \item гетерогенность патчей на границе
          (ср.\ Раздел~\ref{sec:boundary-extended}).
\end{itemize}

Совместная геометрия порогов определяет:
\begin{itemize}
    \item глобальную форму \(\Omega(K)\),
    \item допустимые траектории и структуру циклов,
    \item доступные переходы континуума.
\end{itemize}

\subsection{Размерностные пороги}

Размерностные пороги занимают привилегированное положение в OC, поскольку
они решают, когда новые различия становятся непредставимыми внутри текущих осей.

Пусть \(T(K,t)\) обозначает структурное напряжение.
Тогда, схематично:
\[
  T(K,t) < \Theta_{\rm dim}(K)
   \quad\Rightarrow\quad
   \text{все релевантные различия могут быть выражены в } A(K),
\]
\[
  T(K,t) = \Theta_{\rm dim}(K)
   \quad\Rightarrow\quad
   \text{проекция на } A(K)\ \text{проваливается},
\]
\[
  T(K,t) > \Theta_{\rm dim}(K)
   \quad\Rightarrow\quad
   \text{для структуры требуется новая ось } A_{\rm new}.
\]

Размерностные пороги, в частности, определяют:
\begin{itemize}
    \item рождение \(K_1\) из структурного субстрата \(K_0\),
    \item возникновение химических осей в \(K_3\),
    \item мембранные оси в \(K_4\),
    \item оси возбуждения в \(K_5\),
    \item когнитивные репрезентационные оси в \(K_6\),
    \item социально-нормативные оси в \(K_7\),
    \item мета-логические оси в \(K_9\)–\(K_{10}\).
\end{itemize}

\subsection{Каскады порогов}

Пороги часто взаимодействуют каскадно по форме
\[
  \Theta_{\rm grad}
  \;\longrightarrow\;
  \Theta_{\rm perm}
  \;\longrightarrow\;
  \Theta_{\rm mem},
\]
или более общо: градиенты \(\to\) проницаемость \(\to\) структурный отказ.

Представительные примеры:

\begin{itemize}
    \item \textbf{Взрыв везикулы (\(K_4\)).}  
          Осмотические градиенты превышают \(\Theta_{\rm grad}\), проницаемость растет
          (\(\Theta_{\rm perm}\)), и теряется целостность мембраны
          (\(\Theta_{\rm mem}\)).
    \item \textbf{Потенциалы действия (\(K_5\)).}  
          Деполяризация превышает порог градиента, каналы открываются при
          \(\Theta_{\rm crit}\), и активируется цикл спайка.
    \item \textbf{Институциональный коллапс (\(K_7\)).}  
          Дефицит доверия превышает \(\Theta_{\rm grad}\), нормативная когерентность
          проваливается (\(\Theta_{\rm expr}\)), и институты гибнут
          (\(\Theta_{\rm death}\)).
\end{itemize}

Такие каскады играют центральную роль для ранней жизни, нервных систем и
социальных и инфраструктурных динамик.

\subsection{Пороги смерти}

Пороги смерти определяют момент, когда континуум перестает существовать как живой
континуум. Ниже перечислены эквивалентные проявления смерти в структурном виде:
\begin{itemize}
    \item полное нарушение порогов во всех состояниях:
          \(\forall s:\ \exists k\ \text{с таким, что } f_k(s) > 0\);
    \item исчезновение поддерживающих циклов: \(C(K)=\emptyset\);
    \item коллапс допустимого пространства состояний: \(\Omega(K)=\emptyset\);
    \item репрезентативный отказ, например превышение \(\Theta_{\rm expr}\);
    \item несовместимость с пространством вложения (ни одна конфигурация
          не удовлетворяет одновременно внутренним и вложенным ограничениям).
\end{itemize}

После того как \(\Omega(K)=\emptyset\), смерть структурно необратима:
ни один оператор внутри того же уровня не может воссоздать \(\Omega(K)\);
любой кажущийся возврат соответствует рождению нового континуума \(K'\).

\subsection{Примеры по уровням}

Параметризация порогов проявляется по-разному на каждом уровне иерархии.

\paragraph{\texorpdfstring{\(K_1\)}{K_1} (геометрические континуумы).}
\begin{itemize}
    \item \(\Theta_{\rm exist}\): условия дифференцируемости и регулярности,
          задающие классическую область \(\Omega_{\rm cl}\),
    \item \(\Theta_{\rm crit}\): классические бифуркации и потеря устойчивости.
\end{itemize}

\paragraph{\texorpdfstring{\(K_2\)}{K_2} (физические континуумы).}
\begin{itemize}
    \item \(\Theta_{\rm crit}\): фазовые переходы, включая пороги BKT,
    \item \(\Theta_{\rm dim}\): пороги когерентности, связанные с
          генерацией массы и упорядочиванием,
    \item \(\Theta_{\rm death}\): границы конфайнмента/декуплинга, где
          соответствующий континуум прекращает существование.
\end{itemize}

\paragraph{\texorpdfstring{\(K_3\)}{K_3} (химические континуумы).}
\begin{itemize}
    \item \(\Theta_{\rm exist}\): каталитическое замыкание и существование RAF,
    \item \(\Theta_{\rm crit}\): перколяция реакционной сети,
    \item \(\Theta_{\rm expr}\): недостаточность химических осей для представления
          необходимого разнообразия реакций.
\end{itemize}

\paragraph{\texorpdfstring{\(K_4\)}{K_4} (пребиотические/биологические мембранные континуумы).}
\begin{itemize}
    \item \(\Theta_{\rm grad}\): ограничения осмотического давления,
    \item \(\Theta_{\rm perm}\): пороги проницаемости,
    \item \(\Theta_{\rm mem}\): пороги разрыва мембраны и кривизны.
\end{itemize}

\paragraph{\texorpdfstring{\(K_5\)}{K_5} (возбудимые континуумы).}
\begin{itemize}
    \item \(\Theta_{\rm crit}\): пороги инициации спайка,
    \item \(\Theta_{\rm stab}\): рефрактерные условия и пределы возбудимости,
    \item \(\Theta_{\rm death}\): утрата возбудимости и разрушение
          циклов возбудимых континуумов.
\end{itemize}

\paragraph{\texorpdfstring{\(K_6\)}{K_6} (когнитивные континуумы).}
\begin{itemize}
    \item \(\Theta_{\rm expr}\): пределы связывающей емкости и репрезентационной
          сложности,
    \item \(\Theta_{\rm crit}\): пороги ошибки прогноза,
    \item \(\Theta_{\rm stab}\): пределы когерентности памяти и устойчивости модели.
\end{itemize}

\paragraph{\texorpdfstring{\(K_7\)}{K_7}–\texorpdfstring{\(K_8\)}{K_8} (социальные и цивилизационные континуумы).}
\begin{itemize}
    \item \(\Theta_{\rm exist}\): пороги доверия и легитимности,
    \item \(\Theta_{\rm expr}\): институциональная емкость и репрезентативные пределы,
    \item \(\Theta_{\rm death}\): системный коллапс и разрушение
          институциональных и инфраструктурных циклов.
\end{itemize}

\paragraph{\texorpdfstring{\(K_9\)}{K_9}–\(K_{10}\) (теоретические и мета-теоретические континуумы).}
\begin{itemize}
    \item \(\Theta_{\rm expr}\): выразительная емкость теорий и
          метаязыков,
    \item \(\Theta_{\rm dim}\): пороги расширения мета-уровня,
    \item \(\Theta_{\rm death}\): несогласованность, инконсистентность или структурная
          несовместимость в пространстве моделей.
\end{itemize}

\subsection{Итог}

Расширенный ландшафт порогов дает структурный каркас для всех
качественных изменений в OC. Пороги:
\begin{itemize}
    \item формируют \(\Omega(K)\) и задают \(\partial\Omega(K)\),
    \item управляют фазовыми переходами, возникновением размерности и коллапсом,
    \item опосредуют каскады, связывающие градиенты, проницаемость и структурный
          отказ,
    \item кодируют выразительные и ограничения вложения, способные уничтожать
          континуумы даже без перегрузки энергией.
\end{itemize}

Вместе с геометрией границ и патчей из
Раздела~\ref{sec:boundary-extended} ландшафт порогов завершает
описание того, как континуумы рождаются, поддерживаются и уничтожаются в рамках
OC.


